% from ega3-1-fr.tex
\oldpage[0\textsubscript{III}]{62}
La question étant locale sur $Y$, on peut supposer qu'il existe une suite
exacte
\[
0\longrightarrow\mathcal R\longrightarrow\mathcal O_Y^n
\longrightarrow\mathcal F\longrightarrow0,
\]
et $\mathcal R$ est alors aussi un $\mathcal O_Y$-Module cohérent
(\hyperref[0.5.3.4-fr]{0\textsubscript{I}, 5.3.4}). Pour prouver que les
homomorphismes
\[
f^*(\mathcal{E}xt_{\mathcal O_Y}^p(\mathcal F,\mathcal G))
\longrightarrow
\mathcal{E}xt_{\mathcal O_X}^p(f^*(\mathcal F),f^*(\mathcal G))
\]
sont bijectifs, raisonnons par récurrence sur $p$, la proposition résultant de
(\hyperref[0.6.7.6.1-fr]{0\textsubscript{I}, 6.7.6.1}) lorsque $p=0$. Or, on
a le diagramme commutatif
\[
\hbox{\tiny$
\xymatrix@C=2pt@R=24pt{
f^*(\mathcal{E}xt_{\mathcal O_Y}^{p-1}(\mathcal O_Y^n,\mathcal G))
\ar[r]\ar[d]&
f^*(\mathcal{E}xt_{\mathcal O_Y}^{p-1}(\mathcal R,\mathcal G))
\ar[r]^\partial\ar[d]&
f^*(\mathcal{E}xt_{\mathcal O_Y}^{p}(\mathcal F,\mathcal G))
\ar[r]\ar[d]&
f^*(\mathcal{E}xt_{\mathcal O_Y}^{p}(\mathcal O_Y^n,\mathcal G))
\ar[d]\\
\mathcal{E}xt_{\mathcal O_X}^{p-1}(\mathcal O_X^n,f^*(\mathcal G))
\ar[r]&
\mathcal{E}xt_{\mathcal O_X}^{p-1}(f^*(\mathcal R),f^*(\mathcal G))
\ar[r]^\partial&
\mathcal{E}xt_{\mathcal O_X}^{p}(f^*(\mathcal F),f^*(\mathcal G))
\ar[r]&
\mathcal{E}xt_{\mathcal O_X}^{p}(\mathcal O_X^n,f^*(\mathcal G)).
}
$}
\]
puisque $f^*(\mathcal O_Y)=\mathcal O_X$; comme $f^*$ est exact, les deux
lignes sont exactes. En outre, on a
\[
\mathcal{E}xt_{\mathcal O_Y}^p(\mathcal O_Y^n,\mathcal G)=0
\quad\hbox{pour tout }p>0
\quad\hbox{et de même}\quad
\mathcal{E}xt_{\mathcal O_X}^p(\mathcal O_X^n,f^*(\mathcal G))=0
\quad\hbox{pour tout }p>0
\]
(T, 4.2.3). Vu l'hypothèse de récurrence, les deux premières flèches
verticales du diagramme précédent sont des isomorphismes, et les termes de
droite sont 0, donc
\[
f^*(\mathcal{E}xt_{\mathcal O_Y}^p(\mathcal F,\mathcal G))
\longrightarrow
\mathcal{E}xt_{\mathcal O_X}^p(f^*(\mathcal F),f^*(\mathcal G))
\]
est un isomorphisme.
\end{proof}

\subsection{Hypercohomologie du foncteur image directe}
\label{subsection:0.12.4-fr}

\begin{env}[12.4.1]
\label{0.12.4.1-fr}
Soient $(X,\mathcal O_X)$, $(Y,\mathcal O_Y)$ deux espaces annelés,
$f:X\to Y$ un morphisme d'espaces annelés. On peut prendre
l'\emph{hypercohomologie} de $f_*$ par rapport à un complexe quelconque de
$\mathcal O_X$-Modules
$\mathcal K^\bullet=(\mathcal K^i)_{i\in\mathbb Z}$
(\hyperref[0.11.4.4-fr]{11.4.4}), car dans la catégorie abélienne des
$\mathcal O_Y$-Modules, les limites inductives filtrantes existent et sont
exactes (T, 3.1.1). Les $\mathcal O_Y$-Modules d'hypercohomologie
$R^p f_*(\mathcal K^\bullet)$ se noteront aussi
$\mathcal H^p(f,\mathcal K^\bullet)$ ou
$\mathcal H_f^p(\mathcal K^\bullet)$. Rappelons que
$\mathcal H^\bullet(f,\mathcal K^\bullet)$ est la cohomologie du bicomplexe
de $\mathcal O_Y$-Modules $f_*(\mathcal L^{\bullet\bullet})$, où
$\mathcal L^{\bullet\bullet}$ est une résolution injective de
Cartan-Eilenberg de $\mathcal K^\bullet$ dans la catégorie des
$\mathcal O_X$-Modules~; $\mathcal H^\bullet(f,\mathcal K^\bullet)$ est
l'aboutissement de deux suites spectrales
${}'E(f,\mathcal K^\bullet)$ et ${}''E(f,\mathcal K^\bullet)$ dont les
termes $E_2$ sont donnés par
\begin{equation}
\label{0.12.4.1.1-fr}
\tag{12.4.1.1}
{}'E_2^{pq}=\mathcal H^p(\mathcal H^q(f,\mathcal K^\bullet))
\end{equation}
\begin{equation}
\label{0.12.4.1.2-fr}
\tag{12.4.1.2}
{}''E_2^{pq}=\mathcal H^p(f,\mathcal H^q(\mathcal K^\bullet))
\quad(=R^p f_*(\mathcal H^q(\mathcal K^\bullet))).
\end{equation}
On a, dans ces formules, adopté la notation générale $T(A^\bullet)$ pour
le transformé d'un complexe par un foncteur
(\hyperref[0.11.2.1-fr]{11.2.1}), et on écrit
$\mathcal H^p(f,\mathcal F)$ au lieu de $R^p f_*(\mathcal F)$ pour un
$\mathcal O_X$-Module $\mathcal F$. Rappelons encore que la suite
${}'E(f,\mathcal K^\bullet)$ est toujours \emph{régulière}~; les deux
suites spectrales ${}'E(f,\mathcal K^\bullet)$ et
${}''E(f,\mathcal K^\bullet)$ sont \emph{birégulières} lorsque
$\mathcal K^\bullet$ est limité inférieurement,
